\section{The Two Questions of the Seventeenth Sunday}
\label{sec:opening}

The Seventeenth Sunday after Pentecost, \latin{Dominica decima septima post
Pentecosten}, is a Sunday of the second class in the Proper of Time of the 1962
\work{Missale Romanum}. Its Mass stands on pp.~398--400 of the typical
edition, under the marginal numbers 1602--1611. In 2026 it falls on
20~September. The Mass is said in green, with Gloria and Credo. The calendar's
commemoration of Saints Eustace and Companions on that day gives way to the
Sunday, since a Sunday of the second class admits no commemoration but that
of a feast of the same class (\work{Rubricae generales} 111~b; \work{Rubricae
generales Missalis} 434~b). The Preface is that of the Most Holy Trinity, which
the Missal appoints for every Sunday of the second class outside the Christmas
and Easter seasons. The September Ember Days follow on the Wednesday, Friday
and Saturday of the same week.

At the centre of the Mass stand two questions asked in the Temple at
Jerusalem in the last days before the Passion. A doctor of the law asks Jesus,
``Master, which is the great commandment in the law?'' He is given two
commandments for one: the love of God with the whole heart, the whole soul and
the whole mind, and the love of neighbour, ``like to this''. Then Jesus asks
the Pharisees a question of his own: ``What think you of Christ? Whose son is
he?'' When they answer ``David's'', he presses them with the psalm in which
David, speaking ``in spirit'', calls the Christ his Lord. The Gospel ends in
silence: ``no man was able to answer him a word: neither durst any man from
that day forth ask him any more questions.''

Around these two questions the Missal gathers texts of other books and other
ages. Three of them are prayers of petition in the mouth of a people: the
Introit's plea to be dealt with in mercy and not in justice, Daniel's prayer
for the desolate sanctuary at the Offertory, and the three orations, which ask
to escape the devil's contagion, to be stripped of past and future sins, and to
have our vices cured. Two are praise: the Gradual's blessing of the nation
whose God is the Lord, and the Communion's summons to pay vows to the God who
is terrible to the kings of the earth. One is a cry for a hearing, the
Alleluia's verse from the fifth penitential psalm. And one is an apostle's
appeal from prison, the Epistle, for a unity kept in humility and patience and
grounded in ``one Lord, one faith, one baptism, one God and Father of all''.
Each stands in the table below in the order of the Mass.

The day's texts thus put three questions to those who pray them: what God
commands to be loved, and how; who the Christ is whom David calls Lord; and in
what words a people that knows its own injustice stands before a just God.
The continuation of Guéranger's \work{The Liturgical Year} (Dom Lucien Fromage)
records that the present Gospel gave the day its name, ``the
Sunday of the love of God''. The Introit, the Offertory and the Communion,
which already stand together at this Sunday in the earliest graduals, speak
instead of justice, mercy and vows. These threads, taken in turn, give the
three readings of the formulary that follow. The first hears the whole Mass from the
Gospel's first question, as the prayer of a people learning to love God
wholly and to keep the unity that such love makes. The second hears it from
the Gospel's second question, as the confession of the one God in his Word and
his Spirit, whose Son is David's Lord. The third hears it from the Introit and
the Offertory, as the prayer of a humbled people that confesses God's justice
and receives his mercy. The Fathers and saints who taught about the
particular texts supply all three, every element of the formulary has a place
in each, and four senses of its own, literal, allegorical, moral and
anagogical, close each reading. The three share more than they divide.

\subsection{The appointed elements at a glance}

\begin{maptable}
Introit & \latin{Iustus es, Domine}; Ps 118:137, 124; verse 1 &
God is just and his judgment right; the servant asks to be dealt with
according to God's mercy.\\
Collect & \latin{Da, quaesumus, Domine, populo tuo} &
That the people may shun the devil's contagion and follow God alone with a
pure mind.\\
Epistle & Eph 4:1--6 &
Paul, a prisoner, pleads for a life worthy of the calling: humility,
mildness, patience, mutual support in charity, the unity of the Spirit in the
bond of peace; one body, one Spirit, one hope, one Lord, one faith, one
baptism, one God and Father of all.\\
Gradual & \latin{Beata gens}; Ps 32:12, 6 &
Blessed the nation whose God is the Lord, the people he chose for his
inheritance; the heavens made firm by his word, their power by the spirit of
his mouth.\\
Alleluia & \latin{Domine, exaudi}; Ps 101:2 &
``Hear, O Lord, my prayer: and let my cry come to thee.''\\
Gospel & Mt 22:34--46 &
The great commandment and the second like it; Christ's question about the
Son of David whom David calls Lord; the questioners' silence.\\
Offertory & \latin{Oravi Deum meum ego Daniel}; from Dan 9:17--19 &
Daniel prays that God hear his servant, make his face shine on the sanctuary
and look with favour on the people called by his name.\\
Secret & \latin{Maiestatem tuam, Domine} &
That the holy things we do may strip us of past sins and of future ones.\\
Communion & \latin{Vovete et reddite}; Ps 75:12--13 &
Vow and pay to the Lord, all you round about him who bring gifts, to the God
who takes away the spirit of princes.\\
Postcommunion & \latin{Sanctificationibus tuis} &
That by God's sanctifying gifts our vices may be cured and eternal remedies
come to us.\\
\end{maptable}
